\chapter{\IfLanguageName{dutch}{Stand van zaken}{State of the art}}%
\label{ch:stand-van-zaken}


Waar het vorige hoofdstuk de praktische nood bij Delhaize Ledeberg schetste voor een
camera aan de broodsnijmachine die de klant vertelt welke broodzak hij moet
gebruiken, onderzoekt dit hoofdstuk de technische achtergrond die nodig is om die
nood in te vullen. Het automatisch herkennen van een brood op basis van
een foto kent daarbij twee niveaus. Het praktische minimumdoel is om te bepalen in
welke \emph{broodzakcategorie} een brood thuishoort. Wat een relatief grofmazige
indeling is naargelang de vorm of het onderscheid tussen gewone broden en
biobroden. Een verdergaand en theoretischer doel is om de \emph{specifieke broodsoort}
te onderscheiden, wat gezien het assortiment van ongeveer vijftig soorten een veel fijnmaziger classificatieprobleem vormt.


Om beide niveaus wetenschappelijk te kaderen, bouwt dit hoofdstuk stap voor stap op.
Eerst wordt besproken hoe computervisie vandaag wordt ingezet om brood en andere
voedingsproducten te herkennen, en waarom het onderscheiden van gelijkende
broodsoorten een apart, fijnmazig probleem is. Vervolgens worden de twee mogelijke
taakvormen --- beeldclassificatie en objectdetectie --- tegenover elkaar geplaatst.
Daarna komt transfer learning aan bod als techniek om ook met een beperkte dataset
bruikbare modellen te trainen, gevolgd door de vraag hoe groot de zelf te verzamelen
dataset moet zijn en hoe data-augmentatie daarbij kan helpen.


\section{Computer vision voor brood en voedingsproducten}


Voordat een systeem een brood kan herkennen, moet het bepalen welke visuele kenmerken
het daarvoor gebruikt. \textcite{Olakanmi2023} geven in hun onderzoek een
beeld van hoe computervisie vandaag wordt ingezet om brood en andere bakproducten te
beoordelen. Beeldsystemen baseren zich daarbij op eigenschappen als kleur, textuur,
grootte en vorm om een product te typeren. Diezelfde review wijst er echter op dat
zulke systemen gevoelig zijn voor verstorende factoren zoals belichting, vocht en
ruis~\autocite{Olakanmi2023}. Dit zijn sommige omstandigheden die ook in een winkel voorkomen en waarmee ons systeem zal moeten kunnen omgaan.


Dat computers die visuele kenmerken daadwerkelijk kunnen leren, blijkt uit eerder
broodonderzoek. Zowel \textcite{Cotrim2020} als \textcite{Tantiphanwadi2022} slagen
erin om met convolutionele neurale netwerken (CNN's) de bruiningsgraad van brood te
herkennen, een taak waarbij een brood in enkele duidelijk verschillende stadia wordt
ingedeeld. Dit soort grofmazige classificatie, met een beperkt aantal goed te
onderscheiden klassen, sluit sterk aan bij ons praktisch minimumdoel: ook daar
volstaat het om een brood in een van enkele verschillende
broodzakcategorieën te plaatsen.


Herkenning in een realistische retailcontext is echter een stuk
lastiger. \textcite{Ciapas2023} bouwen een classificatiesysteem voor producten zonder
barcode aan de zelfscankassa (groenten, fruit, snoep,...) en botsen daarbij op typische moeilijkheden: producten
die deels bedekt zijn door de handen van de klant of in zakjes
zitten, en het ontbreken van een bruikbare publieke dataset. In dezelfde lijn
vergelijken \textcite{Boyd2024} meerdere CNN-architecturen voor het herkennen van
afzonderlijke voedingsproducten, en ontwikkelen \textcite{Schmitt2025} met de
Semmeldetector een systeem dat bakproducten (croissants, koeken, brood,...) in een beeld herkent, telt en van elkaar
onderscheidt. Opvallend is dat alle drie de onderzoeken vaststellen dat
bestaande datasets de afzonderlijke producten onvoldoende fijn opsplitsen, wat telkens een eigen dataset noodzakelijk maakt.


Een brood toewijzen aan een brede broodzakcategorie lijkt sterk op de grofmazige taken die in
eerder broodonderzoek al succesvol bleken. De specifieke broodsoort bepalen is
daarentegen verwant aan de moeilijkere, fijnmazige productherkenning, waar klassen
sterk op elkaar kunnen lijken. Dit wordt later verder in detail besproken.


\subsection{Broodsoorten onderscheiden als fijnmazig probleem}


Het onderscheiden van sterk gelijkende broodsoorten sluit nauw aan bij wat in de
literatuur \emph{fine-grained} (fijnmazige) beeldherkenning heet. \textcite{Wei2022}
omschrijven dit als het probleem waarbij de klassen onderling sterk op elkaar lijken, de zogenaamde kleine inter-class-variatie, terwijl de exemplaren binnen
eenzelfde klasse juist sterk kunnen verschillen. Wat men grote intra-class-variatie
noemt. Vertaald naar brood komt dat neer op dit:
Enerzijds kunnen twee verschillende broodsoorten bijna identiek ogen en soms enkel
verschillen in de textuur van de korst of de zaadjes erop. Anderzijds zien twee broden
van exact dezelfde soort er door het bakproces nooit helemaal hetzelfde uit.


Dat verklaart meteen waarom onze twee herkenningsniveaus niet even moeilijk zijn. Een
brood indelen in een brede broodzakcategorie is relatief grofmazig, omdat die
categorieën onderling duidelijker verschillen. De exacte broodsoort bepalen is een
typisch fine-grained probleem, waarbij het model subtiele details moet oppikken om
zeer gelijkende soorten uit elkaar te houden.


\section{Beeldclassificatie versus objectdetectie}


Een tweede fundamentele keuze betreft het soort computervisietaak. \textcite{Liu_2022}
omschrijven beeldclassificatie als het toekennen van één label aan een volledig beeld.
Objectdetectie gaat een stap verder: daar worden objecten in een beeld eerst
\emph{gelokaliseerd}, met een kader errond, en vervolgens geclassificeerd. Een
voorbeeld hiervan is de Semmeldetector van \textcite{Schmitt2025}, die bakproducten
niet alleen herkent maar ze ook in het beeld lokaliseert en telt. Classificatie geeft
dus één antwoord per beeld, terwijl detectie bepaalt \emph{waar} zich \emph{welk}
object bevindt.


In onze opstelling wordt het brood vanuit een vast camerastandpunt aan de
broodsnijmachine gefotografeerd, telkens op een semi-voorspelbare positie en in principe
met één exemplaar per beeld. Onder die voorwaarden voegt de lokalisatie van
objectdetectie weinig toe: er hoeft niets geteld of gesitueerd te worden, enkel het
juiste label bepaald, onafhankelijk of dat label nu een broodzakcategorie of een specifieke
broodsoort is. Voor beide herkenningsniveaus is beeldclassificatie daarom theoretisch
de eenvoudigste en meest directe keuze, ook omdat dit minder calculaties vergt en
minder annotatiewerk vraagt.


Toch worden in dit onderzoek beide taken experimenteel getest. Dit om de oplossingsdeelvraag te beantwoorden over welk model het meest geschikt is. Bovendien zijn reële beelden zelden ideaal: \textcite{Ciapas2023} tonen
dat producten in de praktijk vaak deels bedekt zijn, bijvoorbeeld doordat de klant het
brood vasthoudt. Een detector die het brood eerst lokaliseert en van de achtergrond
scheidt, kan in zulke rommelige omstandigheden alsnog meerwaarde bieden.


Een praktisch voordeel is dat beide sporen grotendeels op dezelfde beelden kunnen
steunen. Dezelfde ruwe foto's dienen zowel voor classificatie als voor detectie; enkel
de annotatie verschilt. Voor classificatie volstaat één label per beeld, terwijl
objectdetectie vereist dat elk brood bijkomend met een \emph{bounding box} wordt
afgebakend. Eén set opnames levert zo twee datasets op, al brengt het detectiespoor
extra annotatiewerk met zich mee.


\section{Transfer learning en beperkte datasets}


Het trainen van een diep neuraal netwerk vanaf nul vereist doorgaans immens grote
hoeveelheden gelabelde data. Transfer learning biedt hiervoor een uitweg.
\textcite{Zhuang2021} beschrijven het als het hergebruiken van kennis uit een
brondomein om het leren in een verwant doeldomein te verbeteren. Concreet betekent dit
dat een netwerk dat al op een grote dataset zoals ImageNet is getraind, en daar
algemene visuele kenmerken zoals randen, texturen en vormen heeft geleerd, als
vertrekpunt dient en vervolgens op de eigen, kleinere dataset wordt bijgesteld.


Welk voorgetraind model het best past, is geen triviale keuze. \textcite{AbouBaker2022}
evalueren elf voorgetrainde ImageNet-netwerken op vijf doeldomeinen en beoordelen ze
niet alleen op nauwkeurigheid, maar ook op trainingstijd en modelgrootte. Zij
benadrukken dat veel studies uitsluitend op nauwkeurigheid focussen en andere, voor de
praktijk belangrijke eigenschappen verwaarlozen~\autocite{AbouBaker2022}. Voor onze
casus is dat zeer relevant: de klant aan de broodsnijmachine mag niet lang
wachten, waardoor de modelkeuze ook door inferentiesnelheid en modelomvang wordt
bepaald en niet enkel door de hoogst behaalde nauwkeurigheid.


Dat transfer learning ook in een broodcontext met weinig data werkt, blijkt opnieuw uit
het werk van \textcite{Tantiphanwadi2022}. Zij combineren een voorgetraind
VGG-16-netwerk met fine-tuning en data-augmentatie, en behalen zo sterke resultaten op
een kleine brooddataset. \textcite{Cotrim2020} komen met compacte CNN's tot een
vergelijkbare vaststelling. Samen bevestigen ze dat zorgvuldig opgezette CNN's in
dit domein geen massale datasets nodig hebben, wat voor een bachelorproef met beperkte
middelen een belangrijke voorwaarde is.


\section{Omvang van de eigen dataset}


Een belangrijke vraag bij het opzetten van dit computervisie-experiment is hoeveel
beelden er precies nodig zijn. Daarop bestaat geen algemeen, vast antwoord: de
benodigde hoeveelheid data hangt af van de taak, de variatie in de beelden, de
gelijkenis tussen de klassen en het gekozen model \autocite{Shorten2019}.


Omdat onze casus specifiek de broodsoorten van Delhaize Ledeberg betreft en er zoals
eerder bleek vermeld, geen publieke dataset bestaat die deze soorten als afzonderlijke
klassen bevat~\autocite{Schmitt2025,Ciapas2023}, stellen we zelf een dataset samen. We
leggen die beelden bewust vast op het niveau van de specifieke broodsoort. Dat is een
doordachte ontwerpkeuze: door elke broodsoort aan de juiste broodzak te koppelen, kan
dezelfde dataset later zowel voor broodsoortclassificatie als voor de afgeleide,  broodzakcategorie worden gebruikt. Een fijnmazig gelabelde dataset laat
zich immers altijd terugbrengen naar een grovere indeling, maar niet omgekeerd.


Hoeveel beelden per soort zinvol zijn, kunnen we het best inschatten door naar
vergelijkbaar onderzoek te kijken. De studies die eerder in dit hoofdstuk aan bod
kwamen, verschillen daarin sterk. \textcite{Tantiphanwadi2022} classificeerden
broodbeelden in drie bruiningsstadia met VGG-16 transfer learning en fine-tuning, en
gebruikten daarvoor 900 originele trainingsbeelden, oftewel 300 per klasse; dit is
brood-specifiek en een classificatietaak, maar met veel minder klassen dan onze
broodsoorten. \textcite{Schmitt2025} onderscheidden met hun Semmeldetector achttien
soorten bakproducten op basis van slechts 1\,151 originele beelden, die ze via een
copy-paste-techniek sterk uitbreidden; dit onderzoek ligt inhoudelijk dicht bij het
onze, maar betreft objectdetectie. \textcite{Boyd2024} onderzochten fijnmazige
voedselherkenning over twintig voedselklassen met meerdere CNN-architecturen, en
beschikten daarvoor over een veel grotere dataset van 41\,949 beelden; net als onze
broodsoorttaak proberen zij sterk gelijkende voedselklassen te onderscheiden, zij het
met aanzienlijk meer data. Tabel~\ref{tab:datasetomvang} zet deze
studies naast elkaar.


\begin{table}
    \centering
    \resizebox{\textwidth}{!}{%
        \begin{tabular}{llccc}
            \toprule
            \textbf{Studie}              & \textbf{Toepassing}                        & \textbf{Klassen} & \textbf{Originele beelden} & \textbf{Beelden per klasse} \\
            \midrule
            \textcite{Tantiphanwadi2022} & bruiningsstadia van brood (classificatie)  & 3                & 900                        & 300 (gebalanceerd)          \\
            \textcite{Schmitt2025}       & bakproducten detecteren/tellen (detectie)  & 18               & 1\,151                     & gem. $\approx$64            \\
            \textcite{Boyd2024}          & fijnmazige voedselherkenning (classificatie) & 20             & 41\,949                    & gem. $\approx$2\,097        \\
            \bottomrule
        \end{tabular}%
    }
    \caption[Datasetomvang in vergelijkbare studies]{\label{tab:datasetomvang}%
        Trainingsdata in vergelijkbare computervisiestudies rond brood- en
        voedselherkenning. De vermelde aantallen betreffen originele, niet-geaugmenteerde
        beelden. ``gem.'' duidt op een rekenkundig gemiddelde, berekend over een dataset die niet noodzakelijk
        gelijk over de klassen verdeeld is.\autocite{Tantiphanwadi2022,Schmitt2025,Boyd2024}.}
\end{table}


De tabel maakt vooral duidelijk dat er geen rechtstreeks overdraagbaar standaardgetal
bestaat: de datasetomvang loopt sterk uiteen naargelang het aantal klassen en de aard
van de taak. Toch tekent zich een bruikbare \emph{indicatieve referentieband} af. Aan
de onderkant tonen \textcite{Schmitt2025} dat zelfs enkele tientallen originele beelden
per klasse kunnen volstaan, al leunen zij daarvoor zwaar op copy-paste-augmentatie en
betreft het objectdetectie. Een brood-specifieke classificatiereferentie ligt hoger:
\textcite{Tantiphanwadi2022} werkten met 300 originele beelden per klasse. Aan de
bovenkant gaat fijnmazige voedselherkenning zoals bij \textcite{Boyd2024} richting
enkele duizenden beelden per klasse, maar met veel grotere middelen dan in een
bachelorproef beschikbaar zijn; dat getal vormt dus geen noodzakelijke bovengrens voor
een proof of concept.


Deze waarden zijn geen universele wetenschappelijke onder- of bovengrens, maar
een referentiekader uit vergelijkbaar onderzoek. De hoogste referentie is voor dit onderzoek niet zomaar
haalbaar. Zouden we de 300 beelden per klasse
van \textcite{Tantiphanwadi2022} aanhouden voor de ongeveer vijftig broodsoorten van
Delhaize Ledeberg, dan kwamen we al op zo'n 15.000 originele foto's, wat veel tijd en rekenwerk zou zijn voor een bachelorproef. De literatuur plaatst de
zinvolle hoeveelheid dus ergens tussen enkele tientallen en enkele honderden originele
beelden per broodsoort, waarbij sterkere augmentatie toelaat om dichter bij de
onderkant te werken. Binnen die haalbare zone wordt de
uiteindelijke hoeveelheid in het methodologiehoofdstuk vastgelegd.


Belangrijker nog dan het aantal originele foto's is de variatie binnen elke klasse. Daarom
fotograferen we per broodsoort meerdere verschillende fysieke exemplaren, in plaats van
vele opnames van één enkel brood. \textcite{Wei2022} wijzen er immers op
dat exemplaren binnen eenzelfde klasse sterk kunnen verschillen, en \textcite{Boyd2024}
benadrukken het belang van die realistische variatie. Een dataset die slechts één
exemplaar per soort bevat, mist die natuurlijke variatie en leidt makkelijker tot
overfitting, waarbij het model de specifieke opnames vanbuiten leert in plaats van de
onderliggende kenmerken van de broodsoort~\autocite{Shorten2019}.


\section{Data-augmentatie}


\textcite{Shorten2019} geven een breed overzicht van data-augmentatie: een verzameling
technieken die een trainingsdataset kunstmatig vergroten en diverser maken om
overfitting tegen te gaan wanneer er weinig data beschikbaar is. Veelgebruikte
technieken zijn geometrische transformaties zoals rotatie, spiegeling, bijsnijden en
translatie, aanpassingen in de kleurruimte zoals helderheid, saturatie en tint, en
\emph{random erasing}, waarbij willekeurige delen van het beeld worden
gewist~\autocite{Shorten2019}.


Verschillende van die technieken sluiten rechtstreeks aan bij onze deelvragen. Variatie
in helderheid en kleur bootst wisselende belichtingsomstandigheden na --- een beperking
die ook \textcite{Olakanmi2023} vermelden --- terwijl random erasing een gedeeltelijke
bedekking simuleert, zoals een hand of een zak die het brood deels aan het zicht
onttrekt~\autocite{Ciapas2023}.


Augmentatie vergroot dus het effectieve aantal trainingsvoorbeelden en maakt het model
robuuster, maar het is geen volwaardige vervanging voor echte variatie tussen fysieke
broden. Afgeleide beelden blijven varianten van bestaande opnames en voegen geen nieuwe,
onafhankelijke observaties toe; augmentatie kan overfitting daarom verminderen, maar niet
volledig oplossen~\autocite{Shorten2019}. In deze opzet dient augmentatie dan ook als
aanvulling: ze wordt later toegepast op de originele
beelden en genereert tijdens het trainen extra variatie, bovenop, en niet in plaats
van, de variatie die we met meerdere fysieke exemplaren vastleggen.